% System Design section — drafted 2026-07-06, ~1.5 pages two-column.
% Verified against code; remaining unbacked claims carry \todo{} markers.
% SAE feature-space visualisation, dashboard, and steering stay in §4 (SAE Support).

\section{System Design}
\label{sec:system}

Orrery's design goal is frictionless~\cite{ren2025embeddingatlaslowfrictioninteractive} movement across the whole exploration loop: from a raw dataset to a navigable constellation, and from a promising SAE feature to its causal effect on the model~\cite{templeton2024scaling}. The platform exposes two workflows over one shared backend: a \emph{corpus workflow} --- embed, project, cluster, label, explore --- following the pattern of Latent Scope, Nomic Atlas, and Embedding Atlas, and an \emph{SAE workflow} (\S\ref{sec:sae}) that treats a sparse autoencoder's feature space as just another collection.

\paragraph{Datasets, collections, storage.}
Orrery separates two concepts. A \textbf{dataset} is a set of documents (text, images, or pre-computed vectors) stored and indexed in DuckDB. A \textbf{collection} is one representation of that dataset: an embedding, its 2D/3D projections, and optional topic extractions and SAE activations. One dataset can therefore have many collections --- the same corpus embedded with different models, prompts, or column combinations --- without duplicating the documents, the main limitation we encountered with ChromaDB-only storage. DuckDB acts as the central orchestrator (documents, JSON metadata, projections as native float arrays, topic tables, sparse SAE activations, full-text index), while ChromaDB\footnote{\url{https://github.com/chroma-core/chroma}} stores only IDs and dense vectors for approximate nearest-neighbour search (Figure~\ref{fig:architecture}). A GraphQL API serves the Next.js frontend; long-running jobs (embedding, topic extraction, activation collection) stream progress over WebSockets and are resumable after interruption.

\paragraph{Ingestion.}
A dataset is previewed from HuggingFace or from local JSON, CSV, or Parquet files. The columns to embed are specified with a free-form template that supports both custom prompts and column concatenation (e.g.\ \texttt{\{word\}: \{definition\}} for WordNet). Eight embedding providers are supported --- local: SentenceTransformers (default), Ollama, FlagEmbedding for BGE~\cite{chen-etal-2024-m3}, and PyTorch for Qwen~\cite{zhang2025qwen3embeddingadvancingtext}; via API key: OpenAI, Cohere, Gemini, and HuggingFace --- including task-type prompts for models that use them, such as EmbeddingGemma~\cite{vera2025embeddinggemmapowerfullightweighttext}. The embedding function is instantiated through a factory, hardware is auto-detected (MPS/CUDA/CPU), and all core features run fully offline with local models; datasets that already contain a vector column skip this step. Projections are then computed with PCA and UMAP~\cite{McInnes2018}, in both 2D and 3D, and stored in DuckDB. Projection is the only mandatory step: the dataset is now visualisable, and everything below is optional.

\paragraph{Topic modelling.}
Topic extraction follows BERTopic~\cite{grootendorst2022bertopicneuraltopicmodeling}: reduce, cluster, extract keywords, label. Since embedding and reduction already exist in the platform, we do not import BERTopic itself but fork only its clustering, c-TF-IDF, and labelling components, reusing its LLM labelling prompt. Run with BERTopic's default recipe (HDBSCAN over a 5-dimensional UMAP), our implementation reproduces its silhouette and $C_v$ scores to the fourth decimal.\todo{Commit the parity-check artifact backing this number.} Where prior platforms fix this pipeline (\S\ref{sec:related}), both steps are configurable: clustering can run on a fresh 5-D UMAP of the raw vectors (default), on the stored 2D/3D visualisation coordinates (fast), or on the L2-normalised original vectors, using HDBSCAN, k-means, Gaussian mixtures, or spectral clustering. Labels come from c-TF-IDF keywords or an LLM and can be regenerated or edited by hand. An optional reduction step merges topics hierarchically toward a target count while preserving the originals as subtopics --- a two-level hierarchy the interface reuses for nested colouring.

\paragraph{Per-document SAE activations.}
Any stored collection can additionally be run through Gemma with Gemma Scope SAEs~\cite{lieberum2024gemmascope} attached (any width, any layer). Following \citet{jiang2025interpretableembeddingssparseautoencoders}, token-level activations are max-pooled into a single sparse vector per document ($\sim$100--150 nonzero features), stored in DuckDB. Because each dimension carries an auto-interpreted feature label, the corpus becomes searchable by feature name; collection is batched, reports progress, and is resumable.

\paragraph{Search and filtering.}
A collection can be searched four ways, with results rendered as a highlighted constellation: (i)~\textbf{semantic} search over the original embedding space, also triggered by clicking any point --- the camera flies to it and its nearest neighbours in the \emph{pre-projection} space light up; (ii)~\textbf{text} search, executed server-side in DuckDB with per-column selection, exact or partial matching, and BM25 word-level search via its FTS extension; (iii)~\textbf{SAE feature search}: querying a feature name (e.g.\ \emph{football}) ranks documents by their max activation over matching features; (iv)~\textbf{prompt highlighting} on SAE collections (\S\ref{sec:sae}). The legend doubles as a filter: categorical fields render as toggleable category lists with counts, numeric fields as histograms with draggable range handles, and auto-detected temporal fields get a dedicated range slider that mutes out-of-range points.

\paragraph{Colour.}
Beyond the standard d3 scales, Orrery ships Crameri's perceptually uniform scientific colour maps~\cite{crameri2018scientific}, which avoid scale distortion and remain readable for colour-blind users~\cite{crameri2020misuse}. The numeric legend is an interactive histogram of the field's distribution with draggable minimum, centre, and maximum handles plus a log option, so a scale skewed by outliers is re-anchored in seconds. For datasets that are themselves about colour, a new \emph{Hilbert} scale samples 512 points from an $8^3$ grid of RGB space and orders them along a 3D Hilbert curve, linearising colour space onto a single locality-preserving strip (used in \S\ref{sec:xkcd}). The active colour configuration round-trips through the URL and can be saved per collection.

\paragraph{Rendering at scale.}
The scatter plots are WebGL-accelerated Plotly.js --- but a vendored fork: on every trace update, Plotly re-reads all $n$ points across all $m$ attributes of every existing trace, an $O(n \cdot m)$ pathology that makes real-time overlays (search highlights, selection markers) impossible at scale. Beyond that fix, we use Plotly only for camera control and point rendering: native tooltips, legend, labels, and animations are disabled and re-implemented on overlay canvases that read the WebGL context directly --- tooltips are positioned from projected coordinates, and fly-to animations manipulate the camera directly. Label placement ports the Embedding Projector's collision-avoidance code~\cite{smilkov2016embeddingprojectorinteractivevisualization}: 3D anchors are projected to screen space each frame and greedily packed on a spatial grid, extended with a priority order (cluster labels first, then similarity-ranked point labels) and rank-dependent opacity. Finally, since no scatter-plot library offers volumetric rendering, the \emph{nebula} mode approximates it: after an unsuccessful attempt with three.js bloom, we render feathered sprites around cluster centroids, spread according to cluster extent, in the style of GalaxyThreeJS,\footnote{\url{https://github.com/pickles976/GalaxyThreeJS}} composited over the plot in screen blend mode. The interface sustains 60\,fps at 250k points with nebulae enabled on 8\,GB of RAM.\todo{State test machine and measurement method.}
